\documentclass[11pt]{amsart}
\usepackage[margin=1.15in]{geometry}
\usepackage[T1]{fontenc}
\usepackage{lmodern,microtype}
\usepackage{amsmath,amssymb,amsthm,mathtools}
\usepackage[colorlinks=true,linkcolor=blue,citecolor=blue,urlcolor=blue]{hyperref}

\newtheorem{theorem}{Theorem}[section]
\newtheorem{lemma}[theorem]{Lemma}
\newtheorem{proposition}[theorem]{Proposition}
\newtheorem{corollary}[theorem]{Corollary}
\newtheorem{conjecture}[theorem]{Conjecture}
\theoremstyle{remark}
\newtheorem{remark}[theorem]{Remark}
\numberwithin{equation}{section}

\newcommand{\lin}{c_1}

\hypersetup{pdftitle={Order polynomials of width-two posets can have negative coefficients},pdfauthor={Brandon Li}}

\title{Order polynomials of width-two posets\\ can have negative coefficients}
\author{Brandon Li}
\address{Great Neck, New York, USA}
\email{brandon.li.gn@gmail.com}
\subjclass[2020]{06A07, 05A15, 52B20}
\keywords{order polynomial, width-two poset, Ehrhart positivity, linear extensions, descents}

\begin{document}

\begin{abstract}
Ferroni, Morales and Panova conjectured that every poset of width two has an
order polynomial with nonnegative coefficients. We disprove this conjecture. For
the $61$-element width-two poset $P$ whose skew shape is the $28\times33$
rectangle with the hook $(7,1,1,1,1)$ removed from a corner, the linear
coefficient of the order polynomial is $-233/872541560$. More generally,
removing the hook $(5,1,1,1,1)$ from an $m\times m$ square yields a negative
linear coefficient for every $m\ge35$, giving infinitely many counterexamples. Within the family of posets obtained
from two chains by adding two relations between them, $61$ elements is the minimum size of a
counterexample.
On the positive side, we prove nonnegativity for three classes of width-two
posets: series-parallel posets, posets covered by two chains one of which has at
most four elements, and posets whose incomparability graph has at most one cycle
in each component. For the first and third classes, we also characterize when
the linear coefficient is positive.
\end{abstract}

\maketitle

\section{Introduction}

Let $P$ be a finite poset with $N$ elements. For a positive integer $t$, let
$\Omega_P(t)$ count the maps $f\colon P\to[t]=\{1,\dots,t\}$ satisfying
$f(x)\le f(y)$ whenever $x<y$ in $P$. Stanley showed that $\Omega_P$ is a
polynomial in $t$ of degree $N$, called the \emph{order polynomial} of
$P$~\cite{Stanley1972,EC1}. It can also be viewed as an Ehrhart polynomial via
$\Omega_P(t)=\operatorname{ehr}_{\mathcal O(P)}(t-1)$, where $\mathcal O(P)$ is
the order polytope of $P$~\cite{Stanley1986}. The coefficients of $\Omega_P$
need not be nonnegative; a basic example is an antichain of four elements with a
minimum adjoined~\cite[Section~8.7]{FMP}.

Ferroni, Morales and Panova~\cite{FMP} proved that the cell poset of every skew
Young diagram has an order polynomial with nonnegative coefficients, and they
conjectured the same property for three further classes, the first of which
consists of posets of width two.

\begin{conjecture}[{\cite[Conjecture~7.1]{FMP}}]\label{conj:FMP}
If $P$ is a poset of width two, then $\Omega_P$ has nonnegative coefficients.
\end{conjecture}

Every width-two poset is covered by two chains $\alpha_1<\dots<\alpha_m$ and
$\beta_1<\dots<\beta_n$, and its incomparable pairs $(\alpha_i,\beta_j)$ form a
skew shape $\lambda/\mu$ inside the $m\times n$
rectangle~\cite[Section~8]{CPP}. Ferroni, Morales and Panova denote this poset
by $R_{\lambda/\mu}$ and formulate Conjecture~\ref{conj:FMP} in those terms.
They verified the conjecture by exhaustive computation when both chains have at
most five elements, as well as for the complement zigzags $\overline Z_n$ with
$n\le200$~\cite[Remark~7.3]{FMP}.

\subsection*{Main results}
Our first result shows that the conjecture is false.

\begin{theorem}\label{thm:main}
Let $P$ be the poset formed by two chains $A_1<\dots<A_{28}$ and
$B_1<\dots<B_{33}$ together with the relations $A_1<B_{27}$ and $A_5<B_{33}$ and
their transitive consequences. Then $P$ has width two, $|P|=61$, and
\[
 [t]\,\Omega_P(t)=-\frac{233}{872541560}.
\]
\end{theorem}

The incomparable pairs of $P$ fill the $28\times33$ rectangle except for a
hook $(7,1,1,1,1)$ at one corner. Up to duality, which leaves the order
polynomial unchanged, $P$ is the poset $R_{\lambda/\mu}$ with
$\lambda=(33^{28})$ and $\mu=(7,1,1,1,1)$. Moreover, the linear coefficient is
the only negative coefficient of $\Omega_P$. Extending the same construction
yields infinitely many counterexamples.

\begin{theorem}\label{thm:family}
For $m\ge5$, let $P_m$ be the width-two poset whose skew shape is the
$m\times m$ square with the hook $(5,1,1,1,1)$ removed from a corner. Then
$[t]\,\Omega_{P_m}(t)<0$ for every $m\ge35$, and
\[
 \lim_{m\to\infty}\,2m(2m-1)\,[t]\,\Omega_{P_m}(t)=-\frac38.
\]
\end{theorem}

Among posets of this two-relation type, $61$ elements is minimal
(Remark~\ref{rem:smallest}); the smallest width-two poset with a negative
coefficient remains unknown.

On the positive side, we establish the conjecture for three subclasses of
width-two posets. Recall that a poset is \emph{series-parallel} if it can be
built from singletons by disjoint unions and ordinal sums, and that the
\emph{incomparability graph} of $P$ has vertex set $P$ with an edge for each
incomparable pair.

\begin{theorem}\label{thm:positive}
Let $P$ be a poset of width at most two. Then $\Omega_P$ has nonnegative
coefficients in each of the following cases:
\begin{enumerate}
 \item $P$ is series-parallel.
 \item $P$ is covered by two disjoint chains, one of which has at most four
 elements.
 \item Every connected component of the incomparability graph of $P$ contains at
 most one cycle.
\end{enumerate}
In cases (1) and (3), if $P$ is nonempty, then $[t]\,\Omega_P(t)>0$ if and only
if the comparability graph of $P$ is connected.
\end{theorem}

In particular, case (3) includes every width-two poset whose incomparability
graph is a forest, and it covers the complement zigzags
of~\cite[Remark~7.3]{FMP}, whose incomparability graphs are paths
(Corollary~\ref{cor:zigzag}).

\subsection*{Methods}
All of our proofs focus on the linear coefficient. Ferroni, Morales and Panova
showed that for any family of posets closed under taking order ideals and
filters, nonnegativity of the linear coefficient implies nonnegativity of every
coefficient~\cite[Theorem~3.4]{FMP}. For a width-two poset, we express the
linear coefficient as an alternating sum over increasing matchings in the board
of incomparable pairs (Section~\ref{sec:prelim}). For the counterexample, this
sum admits a closed-form evaluation via inclusion--exclusion; for the three
positive classes, it is controlled by a representation-theoretic expansion
(Section~\ref{sec:sp}), a finite polynomial certificate
(Section~\ref{sec:short}), or a probabilistic and combinatorial argument
(Section~\ref{sec:sparse}).

All computations in this paper are carried out in exact rational arithmetic, and
code reproducing them is available at
\url{https://github.com/brandonlign/width-two}.

\section{Preliminaries}\label{sec:prelim}

Throughout, $P$ denotes a poset with $N\ge1$ elements, and
$\lin(P)=[t]\,\Omega_P(t)$ denotes the linear coefficient of its order
polynomial.

\subsection{Descents and the linear coefficient}
Fix a natural labeling of $P$. In a linear extension of $P$, a descent is a
position at which the label decreases. Let $h_j$ be the number of linear
extensions with exactly $j$ descents, and set $W_P(z)=\sum_j h_jz^j$. The
standard expansion of the order polynomial~\cite[Section~3.15]{EC1} is
\[
 \Omega_P(t)=\sum_j h_j\binom{t+N-1-j}{N}.
\]
Since $\binom{t+N-1-j}{N}$ has roots $t=j-N+1,\dots,j$, its linear coefficient
for $0\le j\le N-1$ is $(-1)^j/\bigl(N\binom{N-1}{j}\bigr)$. Therefore
\begin{equation}\label{eq:linear}
 \lin(P)=\frac1N\sum_j\frac{(-1)^jh_j}{\binom{N-1}{j}}
 =\int_0^1(1-u)^{N-1}\,W_P\!\left(\frac{-u}{1-u}\right)du,
\end{equation}
where the second equality follows from the beta integral. Writing
$\mathcal L_N(z^j)=(-1)^j/\bigl(N\binom{N-1}{j}\bigr)$ for the linear functional
in \eqref{eq:linear}, we have $\lin(P)=\mathcal L_N(W_P)$.

The order polynomial is multiplicative over disjoint unions, and for any
nonempty poset it vanishes at $t=0$. Hence
\begin{equation}\label{eq:disconnected}
 \lin(P)=0\quad\text{whenever the comparability graph of $P$ is disconnected.}
\end{equation}

\begin{lemma}[{\cite[Theorem~3.4]{FMP}}]\label{lem:reduction}
Let $\mathcal F$ be a family of posets closed under taking order ideals and
filters. If $\lin(P)\ge0$ for every $P\in\mathcal F$, then every $P\in\mathcal F$
has an order polynomial with nonnegative coefficients.
\end{lemma}

\subsection{Width two and increasing matchings}
Suppose $P$ is covered by two disjoint chains $A_1<\dots<A_m$ and
$B_1<\dots<B_k$, with $N=m+k$. For each $i$, let $l_i$ be the number of
$A$-elements below $B_i$, and let $u_i$ be the number of $A$-elements not above
$B_i$. Both sequences are weakly increasing and satisfy $l_i\le u_i$, and $B_i$
is incomparable to $A_x$ if and only if $l_i<x\le u_i$. We call
\[
 H=\{(i,x): l_i<x\le u_i\}
\]
the \emph{board} of $P$. Its cells correspond to the incomparable pairs of $P$,
so $H$ may be viewed as the edge set of the incomparability graph, regarded as a
bipartite graph between the two chains. A set of cells is an
\emph{increasing matching} if, once sorted by first coordinate, both coordinates
strictly increase.

\begin{lemma}\label{lem:matchings}
With the natural labeling given by the earliest linear extension, $h_j$ is the
number of increasing matchings of size $j$ in the board $H$.
\end{lemma}

\begin{proof}
A linear extension of $P$ is uniquely determined by the numbers $x_i$ of
$A$-elements preceding $B_i$, which form a weakly increasing sequence satisfying
$l_i\le x_i\le u_i$. Label $P$ according to the earliest linear extension,
namely the one with $x_i=l_i$ for every $i$. Under this labeling, no descent
enters an $A$-element, and a descent enters $B_i$ if and only if $x_i>l_i$ and
$x_i>x_{i-1}$ (with the convention $x_0=0$). Marking the cell $(i,x_i)$ for each
such $i$ produces a set of cells in $H$ that forms an increasing matching.
Conversely, any increasing matching determines a linear extension: the marked
indices take their marked values, and each unmarked value is recovered greedily
via $x_i=\max(l_i,x_{i-1})$. Monotonicity of the endpoints and of the marked
cells ensures that this sequence is feasible.
\end{proof}

By \eqref{eq:linear} and Lemma~\ref{lem:matchings}, the linear coefficient of a
width-two poset depends only on the number of increasing matchings of each size
in its board. For two disjoint chains of sizes $a$ and $b$, the board is the
full $b\times a$ rectangle, and
\[
 W_{a,b}(z)=\sum_j\binom aj\binom bjz^j,
\]
since an increasing matching of size $j$ is uniquely determined by a choice of
$j$ rows and $j$ columns.

\section{The counterexamples}\label{sec:counterexample}

Fix integers $m,k\ge2$, $2\le r\le m$ and $2\le h\le k$. Let $P(m,k,r,h)$ be the
poset formed by two chains $A_1<\dots<A_m$ and $B_1<\dots<B_k$ together with the
relations
\[
 A_1<B_{k-h+1}\qquad\text{and}\qquad A_r<B_k
\]
and their transitive consequences. The resulting comparable cross pairs are
$A_1<B_j$ for $j>k-h$ and $A_i<B_k$ for $i\le r$. In the board, these relations
remove the last $h$ cells of the first column and the first $r$ cells of the
last row from the full $k\times m$ rectangle, deleting a hook with arm $h-1$ and
leg $r-1$ at one corner. Since $A_m$ and $B_1$ are incomparable, $P(m,k,r,h)$
has width exactly two. In the notation of~\cite{FMP}, $P(28,33,5,7)$ is the
poset of Theorem~\ref{thm:main}, and $P(m,m,5,5)$ is the poset $P_m$ of
Theorem~\ref{thm:family}.

\begin{proposition}\label{prop:W}
The descent polynomial of $P=P(m,k,r,h)$ is
\[
 W_P=W_{m,k}-z\left(\sum_{q=0}^{r-1}W_{k-1,q}+\sum_{q=0}^{h-1}W_{m-1,q}-1\right)
 +z^2\sum_{i=0}^{h-2}\sum_{j=0}^{r-2}W_{i,j}.
\]
\end{proposition}

\begin{proof}
By Lemma~\ref{lem:matchings}, $W_P$ counts increasing matchings in the full
$k\times m$ rectangle that avoid the two deleted strips. Because the cells of
each strip share a row or a column, an increasing matching can use at most one
cell from each strip. In the full rectangle, any matching containing the cell
$(k,q+1)$ of the row strip must have all its remaining cells in the first $k-1$
rows and first $q$ columns; such matchings are therefore counted by
$zW_{k-1,q}$ for $0\le q\le r-1$. Similarly, the matchings containing the cell
$(k-q,1)$ of the column strip are counted by $zW_{m-1,q}$ for $0\le q\le h-1$. A
matching that uses one cell from each strip consists of two distinct cells whose
remaining cells lie strictly between them in both coordinates, yielding the
double sum. Since the shared corner cell is subtracted in both single sums, the
term $-1$ restores it once. Applying inclusion--exclusion over the two strips
gives the stated formula.
\end{proof}

\begin{proof}[Proof of Theorem~\ref{thm:main}]
Apply $\mathcal L_N$ to Proposition~\ref{prop:W} with $(m,k,r,h)=(28,33,5,7)$
and $N=61$. The term $W_{m,k}$ vanishes under $\mathcal L_N$, since it is the
descent polynomial of two disjoint chains and its linear coefficient is zero
by~\eqref{eq:disconnected}. Evaluating the remaining finitely many terms using
$\mathcal L_N(z^j)=(-1)^j/\bigl(N\binom{N-1}{j}\bigr)$ yields
$-233/872541560$.

We verified this value by two further independent computations, both of which
directly count weakly order-preserving maps $P\to[t]$ for $t=0,1,\dots,61$: one
sums over the values assigned to $A_1$ and $A_5$, while the other counts
multichains of order ideals. The resulting counts begin $0,1,975,257280$. Newton
interpolation then recovers the order polynomial exactly, confirming the stated
linear coefficient and showing that all remaining coefficients are positive.
\end{proof}

\begin{proof}[Proof of Theorem~\ref{thm:family}]
Set $m=k$ and $r=h=5$, so that $N=2m$. Summing each strip in
Proposition~\ref{prop:W} via the hockey-stick identity and using
$(j+1)\binom5{j+1}=5\binom4j$, while noting that the coefficient of $z^{j+2}$ in
the double sum is $\binom4{j+1}^2$ for $0\le j\le3$, we simplify to obtain
\[
 N(N-1)\,\lin(P_m)=\frac{10\binom{m-1}{4}-\binom{2m-6}{4}+4\binom{2m-6}{3}}
 {\binom{2m-2}{4}}.
\]
Substituting $m=35+x$, the numerator becomes
\[
 -\frac{x^4}{4}-\frac{143x^3}{6}-\frac{3049x^2}{4}-\frac{25001x}{3}-4960,
\]
which is negative for every $x\ge0$. Finally, since the numerator is
$-\tfrac{1}{4}m^4+O(m^3)$ and the denominator is $\tfrac{2}{3}m^4+O(m^3)$, the
displayed formula yields the claimed limit.
\end{proof}

The same formula is positive for $5\le m\le34$, so $m\ge35$ is precisely the
range of counterexamples within this one-parameter family.

\begin{remark}\label{rem:smallest}
Because Proposition~\ref{prop:W} expresses $\lin\bigl(P(m,k,r,h)\bigr)$ as a
short exact sum, the entire family can be searched exhaustively. Every member
with $m+k\le60$ satisfies $\lin\ge0$. For $m+k=61$, the negative members are
$P(28,33,5,7)$ and $P(30,31,6,6)$, with linear coefficients $-233/872541560$ and
$-4563/6107790920$, together with their duals $P(33,28,7,5)$ and
$P(31,30,6,6)$; in general, the dual of $P(m,k,r,h)$ is $P(k,m,h,r)$.
\end{remark}

\section{Series-parallel posets}\label{sec:sp}

\begin{theorem}\label{thm:sp}
Every series-parallel poset of width at most two has an order polynomial with
nonnegative coefficients. If such a poset is nonempty, its linear coefficient is
positive exactly when it is connected.
\end{theorem}

\begin{proof}
Every nontrivial parallel composition of width at most two consists of two
chains. Consequently, any series-parallel poset of width at most two is an
ordinal sum of blocks $C_a\sqcup C_b$, where $C_a$ denotes an $a$-element chain
and either chain may be empty. This class is closed under taking ideals and
filters. The descent polynomial of each block is $W_{a,b}$, and under a
blockwise natural labeling, the descent polynomial of an ordinal sum is the
product of the descent polynomials of its blocks.

Set $\phi_{a,b}(u)=(1-u)^{(a+b)/2}W_{a,b}\bigl(-u/(1-u)\bigr)$. In
$\operatorname{Sym}^{a+b}(\mathbb C^2)$, this function is the diagonal matrix
coefficient, at the unit monomial of bidegree $(a,b)$, of the rotation with
diagonal entries $\sqrt{1-u}$ and off-diagonal entries $\pm\sqrt u$; this
follows by expanding the images of the two variables and extracting the original
monomial. The product over all blocks is therefore the diagonal matrix
coefficient at the unit tensor of these monomials. Decomposing the tensor
product orthogonally into irreducibles $\operatorname{Sym}^d(\mathbb C^2)$ with
$d\le N$ and $d\equiv N\pmod2$, we observe that in each summand the projected
tensor lies in a single weight space, which is at most one-dimensional. Hence
the product is a nonnegative linear combination of functions $\phi_{A,B}$ with
$A+B=d$, weighted by squared projection norms.

By~\eqref{eq:linear}, it remains to integrate $(1-u)^{N/2-1}\phi_{A,B}(u)$ over
$[0,1]$. Setting $M=(N+d)/2-1$, this integral equals
\[
 \frac1{M+1}\sum_j\frac{(-1)^j\binom Aj\binom Bj}{\binom Mj}.
\]
If $A=0$ or $B=0$, the sum equals $1$. Otherwise $M\ge d-1\ge\max(A,B)$, and by
inclusion--exclusion the sum is the probability that a uniformly random
$A$-subset of an $M$-set avoids a fixed $B$-subset. Each summand in the
expansion therefore makes a nonnegative contribution, giving $\lin(P)\ge0$.
Nonnegativity of all coefficients now follows from Lemma~\ref{lem:reduction}.

For strictness, discard any empty factors and set $A_0=\sum_ia_i$,
$B_0=\sum_ib_i$ and $n_i=a_i+b_i$. If $A_0=0$ or $B_0=0$, the poset is a chain.
Otherwise, the top summand $d=N$ has squared projection norm
\[
 \frac{\prod_i\binom{n_i}{a_i}}{\binom N{A_0}}.
\]
Indeed, the normalized completely symmetric weight vector is the sum over all
allocations of the $A_0$ first-variable positions, and a fixed allocation across
the blocks has the stated squared norm. When there are at least two nonempty
blocks and $0<A_0<N$, another allocation exists: one can move a first-variable
position from an occupied block to a different block containing a
second-variable position. Thus the top squared norm is strictly less than one,
so some lower summand has nonzero weight and contributes a positive amount,
since $M\ge A+B$ holds there. Finally, a single block with $a,b>0$ is
disconnected and has linear coefficient zero by~\eqref{eq:disconnected}.
\end{proof}

\section{A short chain}\label{sec:short}

\begin{theorem}\label{thm:short}
If $P$ is covered by two disjoint chains and one of them has at most four
elements, then $\Omega_P$ has nonnegative coefficients.
\end{theorem}

\begin{proof}
Let the long chain have $m$ elements and the short chain $k\le4$ elements. The
case of a single chain is immediate, so assume $k\ge1$. By
Lemma~\ref{lem:matchings}, $h_j$ counts choices of $j$ increasing indices $i$ in
the short chain together with strictly increasing marks $x_i\in(l_i,u_i]$.

Sort the $2k$ endpoints $l_1,\dots,l_k,u_1,\dots,u_k$, breaking ties
consistently with the two orders and with $l_i\le u_i$; this partitions $[0,m]$
into $2k+1$ nonnegative integer gaps. Within each ordering of the endpoints (a
chamber), $h_j$ is a polynomial of total degree at most $j$ in the gap sizes,
obtained by assigning each mark to a gap and choosing the ordered marks within
each gap via a binomial coefficient in its size. Consequently,
\[
 (m+1)N(N-1)_{\underline k}\,\lin(P)
 =(m+1)\sum_j(-1)^j j!\,h_j\,(N-1-j)_{\underline{k-j}}
\]
is a polynomial of degree at most $k+1$ in the gap sizes, where
$(x)_{\underline j}$ denotes the falling factorial. For every chamber, we
computed its coefficients in the multivariate binomial basis by exact integer
binomial inversion on the gap vectors of total size at most $k+1$, and verified
that all of them are nonnegative for $1\le k\le4$. The prefactor is positive
whenever $m\ge1$, while at $m=0$ the right-hand side gives the polynomial
extension with no marks. Thus $\lin(P)\ge0$.

Since any ideal or filter of $P$ is again covered by two chains one of which has
at most four elements, Lemma~\ref{lem:reduction} applies.
\end{proof}

\section{Sparse incomparability graphs}\label{sec:sparse}

We call a set of cells of the board \emph{valid} if it is an increasing
matching. Note that validity is hereditary: every subset of a valid set is
valid.

\begin{theorem}\label{thm:sparse}
Let $P$ be a nonempty poset of width at most two with $N$ elements and at most
$N-1$ incomparable pairs. Then $\lin(P)\ge0$, with equality exactly when the
comparability graph of $P$ is disconnected. In particular, if the
incomparability graph of $P$ is a forest, then $\Omega_P$ has nonnegative
coefficients.
\end{theorem}

\begin{proof}
Pad the board with dummy cells so that it has exactly $N-1$ cells, and declare a
set of cells to be valid only if all of its cells are genuine and form an
increasing matching. Order the padded board uniformly at random, and let $L$ be
the length of the longest valid prefix. Because the prefix of size $j$ is a
uniformly random $j$-subset, Lemma~\ref{lem:matchings} gives
$\Pr(L\ge j)=h_j/\binom{N-1}{j}$. By \eqref{eq:linear},
\[
 N\lin(P)=\sum_{j=0}^{N-1}(-1)^j\Pr(L\ge j)=\Pr(L\text{ is even}).
\]
This probability is positive whenever there is at least one dummy cell or
$N=1$, since in either case $L=0$ occurs with positive probability.

Suppose now that there are exactly $N-1$ genuine cells and $N>1$. Assume first
that no valid pair exists. Then the bipartite incomparability graph contains no
two disjoint edges: two disjoint cells in increasing position already form a
valid pair, while two crossing cells force the other two corners of their
rectangle to belong to the board by monotonicity of the endpoints, and those two
corners form a valid pair. A bipartite graph with at least one edge and no two
disjoint edges consists of a star together with isolated vertices. Having $N-1$
edges, this star must span $P$, so its center is isolated in the comparability
graph.

Otherwise, there exists a valid pair. There must also be an invalid pair, since
$N-1$ edges on $N\ge3$ vertices cannot be pairwise disjoint. Suppose for
contradiction that every third cell extends every valid pair to a valid triple.
Starting from a single valid pair, its first cell then forms a valid pair with
every other cell; applying the same assumption to those pairs implies that every
pair is valid, a contradiction. Hence some valid pair followed by a third cell
that does not extend it yields $L=2$ with positive probability, so $\lin(P)>0$.
Conversely, if the comparability graph is disconnected, then $\lin(P)=0$
by~\eqref{eq:disconnected}.

Finally, the property of being a forest is inherited by induced subgraphs, and
the incomparability graph of an ideal or filter is an induced subgraph of that
of $P$. Therefore the class of width-two posets whose incomparability graph is a
forest is closed under ideals and filters, and Lemma~\ref{lem:reduction}
applies.
\end{proof}

\begin{theorem}\label{thm:pseudoforest}
Let $P$ be a poset of width at most two such that every connected component of
its incomparability graph contains at most one cycle. Then $\Omega_P$ has
nonnegative coefficients. If $P$ is nonempty, then $\lin(P)>0$ if and only if
the comparability graph of $P$ is connected.
\end{theorem}

\begin{proof}
Let $N=|P|\ge1$ and let $H$ be the incomparability graph of $P$. Any graph in
which every component contains at most one cycle has at most $N$ edges. If $H$
has at most $N-1$ edges, then Theorem~\ref{thm:sparse} applies; we may therefore
assume that $H$ has exactly $N$ edges, which occurs only when every component is
unicyclic.

Call a pair of cells in distinct rows and distinct columns a
\emph{crossing pair} if their column indices appear in reverse order. By
monotonicity of the endpoints, a crossing pair forces the other two corners of
its rectangle to be cells of $H$, so the two edges lie on a $4$-cycle.
Conversely, every cycle contains a crossing pair: in the least row of the cycle,
pair the edge at the larger of its two cycle columns with the edge at the
smaller column in another row of the cycle. Consequently, each unicyclic
component has a unique cycle, which has length four, and no leaf edge of $H$
belongs to a crossing pair.

Orient each cycle cyclically and orient every remaining edge toward the cycle of
its component, so that every vertex has out-degree one. For each edge in a valid
$j$-set, the outgoing edge at its head is incompatible with the set. These $j$
incompatible edges are distinct and unselected, since a valid set is a matching
and no edge is oriented in both directions. Counting valid one-edge extensions
therefore yields
\[
 (j+1)h_{j+1}\le(N-2j)h_j\qquad(1\le j\le\lfloor N/2\rfloor).
\]
It follows that $a_j=h_j/\binom{N-1}{j}$ is nonincreasing for $1\le j\le N-1$,
with $a_j=0$ for $j>\lfloor N/2\rfloor$. Since $a_0-a_1=-1/(N-1)$, the
alternating sum in \eqref{eq:linear} shows that it suffices to prove
$a_2-a_3\ge1/(N-1)$, with strict inequality implying $\lin(P)>0$.

For each valid pair $D$, let $f(D)$ denote the number of unselected edges
incompatible with $D$. Because each valid triple extends three distinct valid
pairs, the desired inequality is equivalent to
\[
 S:=\sum_{D\text{ valid},\ |D|=2}\bigl(f(D)-1\bigr)=(N-3)h_2-3h_3\ \ge\ \binom{N-2}{2}.
\]
We establish this bound by iteratively pruning leaves while retaining the
cycles; since pruning a leaf corresponds to an induced restriction of $P$,
compatibility among the remaining edges is unaffected.

Suppose first that $H$ is a disjoint union of $r$ four-cycles. The unique valid
pair inside each cycle has $f(D)=2$. Any two edges from different components
form a valid pair, as they neither cross nor share a vertex, and every such pair
satisfies $f(D)\ge4$ because each edge has at least two incompatible neighbors
in its own cycle. Hence
\[
 S\ge r+48\binom r2=\binom{4r-2}{2}+(r-1)(16r+3).
\]
Now restore one pruned leaf edge $v$, and let $N$ denote the resulting number of
edges. Let $a$ and $b$ be the numbers of old edges incompatible and compatible
with $v$, respectively, so that $a+b=N-1$. We have $a\ge1$ because $v$ meets its
unicyclic component, and $b\ge1$ because some old cycle edge avoids the
attachment vertex and $v$ cannot cross it. Restoring $v$ does not decrease the
contributions of old pairs, while each new valid pair containing $v$ satisfies
$f(D)\ge a$ by definition and $f(D)\ge2$ by the orientation argument. Therefore
$S$ increases by at least
\[
 b\max\{1,a-1\}\ \ge\ a+b-2=N-3=\binom{N-2}{2}-\binom{N-3}{2}.
\]
This inequality is immediate for $a=1$, and for $a\ge2$ the difference between
the two sides is $(a-2)(b-1)\ge0$. This establishes the bound.

The bound is already strict at the base step whenever $r\ge2$. When $r=1$ and
$H$ has at least one leaf, the first restoration makes it strict: the unique old
valid pair of the four-cycle spans all four of its vertices, so its value of
$f(D)$ increases by one in addition to the contribution of the new pairs, and
subsequent restorations preserve strictness. The sole remaining equality case is
a single four-cycle, whose comparability graph is disconnected and which
therefore satisfies $\lin(P)=0$ by~\eqref{eq:disconnected}.

Since induced restrictions preserve both having width at most two and the cycle
condition on $H$, this class is closed under ideals and filters, and
Lemma~\ref{lem:reduction} completes the proof.
\end{proof}

\begin{proof}[Proof of Theorem~\ref{thm:positive}]
Cases (1), (2) and (3) follow directly from Theorems~\ref{thm:sp},
\ref{thm:short} and~\ref{thm:pseudoforest}, respectively.
\end{proof}

\begin{corollary}\label{cor:zigzag}
For every $N$, the complement zigzag $\overline Z_N$ has an order polynomial with
nonnegative coefficients. For $N\ge4$, its coefficients in degrees
$1,\dots,N$ are all positive.
\end{corollary}

\begin{proof}
The complement zigzag $\overline Z_N$ of~\cite[Remark~7.3]{FMP} is the poset on
$[N]$ in which $i<j$ if and only if $j-i\ge2$; the defining cover relations of
lengths two and three generate all such differences. Its incomparability graph
is a path, so Theorem~\ref{thm:sparse} yields nonnegativity, and for $N\ge4$ its
linear coefficient is positive because the comparability graph is connected.
When $N=4$,
\[
 \Omega_{\overline Z_4}(t)=\frac{5t^4+10t^3+7t^2+2t}{24}.
\]
For $N>4$, consider the coefficient convolution arising from the cut recursion
in the proof of~\cite[Theorem~3.4]{FMP}. It includes a term corresponding to the
ideal $\{1\}$ and its complementary filter $\overline Z_{N-1}$. By induction,
this product contributes a positive amount in each degree $2,\dots,N$, while all
remaining contributions are nonnegative.
\end{proof}

\begin{remark}
The proof of Theorem~\ref{thm:sp} expresses the descent polynomial as a
nonnegative linear combination of rectangle-type polynomials of fixed weight.
Such an expansion does not exist for all width-two posets: for two chains of
length three with $l=(0,0,0)$ and $u=(2,3,3)$, direct counting gives
$W_P=1+8z+9z^2+z^3$, whose coordinates in the ordered basis
$(1-z)^rW_{3-r,3-r}(z)$, $0\le r\le3$, are $(57,-5,3,5)/60$. This rules out
that method of proof for general width-two posets, though not nonnegativity
itself.
\end{remark}

\section*{Statements and Declarations}

\paragraph{Use of AI tools.} Claude (Anthropic) helped with the proofs, the verification code, and the
manuscript. The author takes full responsibility for the content.

\paragraph{Competing interests.} The author has no relevant financial or non-financial
interests to disclose.

\paragraph{Funding.} No funding was received for this work.

\paragraph{Data availability.} The verification code is provided as Online
Resource~1 and at \url{https://github.com/brandonlign/width-two}. No other data
were used.

\begin{thebibliography}{9}

\bibitem{CPP}
S.~H. Chan, I.~Pak, and G.~Panova,
\emph{The cross-product conjecture for width two posets},
Trans. Amer. Math. Soc. \textbf{375} (2022), 5923--5961.

\bibitem{FMP}
L.~Ferroni, A.~H. Morales, and G.~Panova,
\emph{Skew shapes, Ehrhart positivity and beyond},
Proc. Lond. Math. Soc. \textbf{133} (2026), e70204.
\href{https://arxiv.org/abs/2503.16403}{arXiv:2503.16403}.

\bibitem{Stanley1972}
R.~P. Stanley,
\emph{Ordered structures and partitions},
Mem. Amer. Math. Soc. \textbf{119} (1972).

\bibitem{Stanley1986}
R.~P. Stanley,
\emph{Two poset polytopes},
Discrete Comput. Geom. \textbf{1} (1986), 9--23.

\bibitem{EC1}
R.~P. Stanley,
\emph{Enumerative Combinatorics, Volume 1}, 2nd ed.,
Cambridge Stud. Adv. Math. \textbf{49}, Cambridge University Press, 2012.

\end{thebibliography}

\end{document}
